% Generated draft: requires mathematical and source review.
\documentclass[11pt]{article}
\usepackage[margin=25mm]{geometry}
\usepackage{fontspec}
\setmainfont{Latin Modern Roman}
\usepackage{amsmath,amssymb,mathtools,unicode-math}
\setmathfont{Latin Modern Math}
\usepackage{enumitem,needspace,xurl,hyperref}
\setlength{\parindent}{0pt}
\setlength{\parskip}{5pt}
\setlength{\emergencystretch}{3em}
\newcommand{\N}{\mathbb{N}}
\newcommand{\Z}{\mathbb{Z}}
\newcommand{\Q}{\mathbb{Q}}
\newcommand{\R}{\mathbb{R}}
\newcommand{\C}{\mathbb{C}}
\newcommand{\F}{\mathbb{F}}
\newcommand{\A}{\mathbb{A}}
\newcommand{\PP}{\mathbb{P}}
\newcommand{\E}{\mathbb{E}}
\newcommand{\eps}{\varepsilon}
\newcommand{\dd}{\,\mathrm{d}}
\begin{document}
\section*{UnsolvedMath ID: 30001918}
\textbf{Source problem code:} catalogue entry
\subsection{Definitions and mathematical statement}
A polarized manifold \((X,L)\) is understood here in the usual projective sense: \(X\) is a smooth complex projective manifold and \(L\) is an ample holomorphic line bundle on \(X\). Write \(c_1(L)\in H^2(X,\R)\) for its real first Chern class. A constant-scalar-curvature K\"ahler metric in this class is a K\"ahler form \(\omega\) satisfying
\[
[\omega]=c_1(L),\qquad \operatorname{Scal}(\omega)\equiv s
\quad\text{for some }s\in\R.
\]
Here the scalar curvature is that of the Riemannian metric determined by \(\omega\) and the complex structure.

The stability terminology requires a qualification. A test configuration of exponent \(r\geq 1\) for \((X,L)\) is a flat projective family \(\pi:\mathcal X\to\A^1_{\C}\), equipped with a \(\C^*\)-action covering the standard action on the base and a relatively ample equivariant line bundle \(\mathcal L\), whose restriction over \(\A^1_{\C}\setminus\{0\}\) is equivariantly isomorphic to the product family with polarization \(L^{\otimes r}\). Its Donaldson--Futaki invariant is denoted by \(\operatorname{DF}(\mathcal X,\mathcal L)\), with the sign convention in which stability requires positivity. Under a usual strict convention, K-stability requires this invariant to be positive for every nontrivial admissible test configuration. K-polystability instead requires nonnegativity, with equality permitted precisely for product configurations, including products induced by nontrivial \(\C^*\)-actions on \((X,L)\). A trivial configuration has trivial action on \(X\), allowing a change of linearization on the line bundle. The admissible class of configurations, including normality and treatment of configurations trivial away from subsets of codimension at least two, must also be specified to obtain an exact definition.

\textbf{Source question (Yau--Tian--Donaldson conjecture, as supplied).} For every polarized manifold \((X,L)\), is the following equivalence true?
\[
\left(\exists\,\omega\text{ K\"ahler on }X:\ [\omega]=c_1(L),\ \operatorname{Scal}(\omega)\text{ is constant}\right)
\quad\Longleftrightarrow\quad
(X,L)\text{ is K-stable}.
\]
\textbf{Unresolved qualification.} The supplied record does not define its use of ``K-stable,'' specify the admissible test configurations, or impose a discreteness hypothesis on automorphisms. In particular, it does not resolve whether K-polystability was intended. The displayed question preserves the source wording; it does not silently replace that wording by K-polystability, uniform stability, or a relative stability condition. These choices must be resolved before this can serve as a fully specified conjecture.
\subsection{Short English statement}
The supplied question asks whether a polarized complex projective manifold admits a constant-scalar-curvature K\"ahler metric in its polarization class exactly when it is K-stable. The meaning of stability needs clarification, especially the distinction between strict K-stability and K-polystability when automorphisms are present. This is a statement-only draft requiring mathematical review.
\subsection{Sources}
\begin{itemize}
\item Supplied local UnsolvedMath registry record, permanent ID \(30001918\), exact source problem code \texttt{catalogue entry}, titled ``The Yau--Tian--Donaldson Correspondence for Polarized Manifolds,'' in the Oberwolfach Reports Open Problems collection. The record attributes the question to the Yau--Tian--Donaldson conjecture; no individual proposer is supplied.
\item Supplied report citation: \emph{Complex Algebraic Geometry (2012)}, Oberwolfach Reports, pp.\ 2669--2733, DOI \texttt{10.4171/owr/2011/47}. \url{https://doi.org/10.4171/owr/2011/47}. The background describes Oberwolfach Reports as documenting mathematical workshops and their open-problem discussions. The registry separately records proposed year \(2011\).
\item Supplied literature reference: Chen--Donaldson--Sun, ``Kahler--Einstein metrics and stability.'' \url{https://arxiv.org/abs/1210.7494}.
\item Supplied literature reference: Chen--Cheng, ``cscK existence results.'' \url{https://arxiv.org/abs/1801.00656}.
\item Supplied literature reference: ``original OWR report.'' \url{https://doi.org/10.4171/owr/2011/47}.
\item Literature assessment dated exactly \texttt{2026-08-21}: the supplied curation labels the problem ``partially solved'' and states that the Yau--Tian--Donaldson correspondence is proved in major settings, including smooth Fano manifolds, that cscK existence is characterized analytically in important cases, and that the full polarized-manifold statement in the excerpt remains open. This is supplied dated curation, not independently verified here. The background explicitly requests independent verification of time-sensitive status and cited claims.
\item The supplied registry labels the statement \texttt{corrected\_verified} and reports that its mathematical content was checked against the official Oberwolfach report, with notation, context, and question form normalized without changing the source-level claim. This is a supplied verification claim; no source has been independently visited or checked for this draft.
\end{itemize}
\end{document}
